\documentclass[11pt]{report}

\usepackage[margin=1in]{geometry}
\usepackage{amsmath,amssymb,mathtools,amsthm}
\usepackage{esint}
\usepackage[most]{tcolorbox}
\usepackage{mathrsfs}
\usepackage[T1]{fontenc}
\usepackage[utf8]{inputenc}
\usepackage{enumitem}
\usepackage{microtype}
\usepackage{hyperref}
\usepackage{titlesec}

\titleformat{\chapter}[display]
  {\normalfont\bfseries\Huge}
  {\chaptername\ \thechapter}
  {0.5em}
  {}

\titlespacing*{\chapter}{0pt}{-30pt}{20pt}

\setlength{\parindent}{0pt}
\setlength{\parskip}{0.6em}
\setlist[enumerate]{itemsep=0.25em, topsep=0.25em}
\setcounter{tocdepth}{2}

\theoremstyle{definition}
\newtheorem{definition}{Definition}[chapter]
\newtheorem{example}[definition]{Example}
\newtheorem{note}[definition]{Note}
\newtheorem{intuition}[definition]{Intuition}

\theoremstyle{plain}
\newtheorem{theorem}[definition]{Theorem}

\theoremstyle{remark}
\newtheorem{remark}[definition]{Remark}

\newcommand{\R}{\mathbb{R}}
\newcommand{\E}{\mathbb{E}}
\newcommand{\T}{\hat T}
\newcommand{\ip}[2]{\left(#1,#2\right)}
\newcommand{\grad}{\nabla}
\newcommand{\divergence}{\operatorname{div}}
\newcommand{\curl}{\operatorname{curl}}
\newcommand{\supp}{\operatorname{supp}}
\newcommand{\sgn}{\operatorname{sgn}}
\newcommand{\Span}{\operatorname{span}}
\newcommand{\ut}{\mathbf{t}}

\begin{document}

\begin{titlepage}
    \centering
    \vspace*{2.8cm}

    {\Huge\bfseries Vector Analysis and Geometry\par}
    \vspace{0.25cm}
    {\huge 2MBC40\par}
    
    \vspace{0.9cm}


    \vspace{2.2cm}



    \vfill

    \centering
    The source code for these notes is publicly available in my github.\\ Contributions and corrections are always welcome.

    \vspace{0.5em}
    \texttt{https://github.com/Luiser13/vector-analysis-summary}\\
    

    \vspace{0.5em}
    {\normalsize Luis Rodriguez Pacheco\par}
    {\normalsize Eindhoven University of Technology\par}
\end{titlepage}

\tableofcontents
\clearpage

\chapter*{Useful Tools Not Covered in the Course}

\begin{center}

\begin{minipage}{0.48\textwidth}
\textbf{Coordinate transformations}

\begin{tcolorbox}[width=\textwidth, colframe=black, colback=white]
\textbf{Polar}
\begin{gather*}
x = r\cos\theta,\quad y = r\sin\theta \\
r\in[0,\infty),\ \theta\in[0,2\pi] \\
\det(D\Phi_{\mathrm{pol}})=r,\quad dA=r\,dr\,d\theta
\end{gather*}
\end{tcolorbox}

\vspace{0.3em}

\begin{tcolorbox}[width=\textwidth, colframe=black, colback=white]
\textbf{Elliptic Polar}
\begin{gather*}
x = x_0 + ar\cos\theta,\quad y = y_0 + br\sin\theta \\
r\in[0,1],\ \theta\in[0,2\pi] \\
\det(D\Phi_{\mathrm{ell}})=abr,\quad dA=abr\,dr\,d\theta
\end{gather*}
\end{tcolorbox}

\vspace{0.3em}

\begin{tcolorbox}[width=\textwidth, colframe=black, colback=white]
\textbf{Cylindrical}
\begin{gather*}
x = r\cos\theta,\quad y = r\sin\theta,\quad z=z \\
r\in[0,\infty),\ \theta\in[0,2\pi],\ z\in\R \\
\det(D\Phi_{\mathrm{cyl}})=r,\quad dV=r\,dr\,d\theta\,dz
\end{gather*}
\end{tcolorbox}

\vspace{0.3em}

\begin{tcolorbox}[width=\textwidth, colframe=black, colback=white]
\textbf{Spherical}
\begin{gather*}
x = \rho\sin\phi\cos\theta \\
y = \rho\sin\phi\sin\theta \\
z = \rho\cos\phi \\
\rho\in[0,\infty),\ \theta\in[0,2\pi],\ \phi\in[0,\pi] \\
\det(D\Phi_{\mathrm{sph}})=\rho^2\sin\phi,\quad\\
dV=\rho^2\sin\phi\,d\rho\,d\phi\,d\theta
\end{gather*}
\end{tcolorbox}

\end{minipage}
\hfill
\begin{minipage}{0.48\textwidth}
\textbf{Weierstrass substitution}

\begin{tcolorbox}[colframe=black, colback=white]
Useful for turning trigonometric integrals into rational integrals.

\begin{gather*}
t=\tan\left(\frac{x}{2}\right),\qquad
dx=\frac{2}{1+t^2}\,dt \\
\sin x=\frac{2t}{1+t^2},\qquad
\cos x=\frac{1-t^2}{1+t^2}
\end{gather*}

\textbf{When to use:}
\begin{itemize}[leftmargin=1.5em]
    \item Integrands involving only \(\sin x\), \(\cos x\), and rational operations.
    \item Especially useful when standard identities do not simplify the integral enough.
\end{itemize}

\textbf{Important:}
For a full period \(x\in[0,2\pi]\), split at \(x=\pi\), since
\[
\tan\left(\frac{\pi}{2}\right) \text{ is undefined.}
\]

\textbf{Procedure:}
\begin{enumerate}[leftmargin=1.5em]
    \item Set \(t=\tan(x/2)\).
    \item Replace \(\sin x,\cos x,dx\) using the formulas above.
    \item Simplify to a rational function in \(t\).
    \item Integrate using algebraic methods: partial fractions, symmetry, completing the square.
\end{enumerate}
\end{tcolorbox}

\end{minipage}

\end{center}

\part{Classical Vector Analysis}

\chapter{Euclidean Space and Fields}

\section{Euclidean Space \texorpdfstring{$\mathbb{R}^3$}{R3}}

\begin{definition}[Inner Product]
\[
\ip{\overrightarrow{OX}}{\overrightarrow{OY}}
=
|\overrightarrow{OX}|\,|\overrightarrow{OY}|
\cos\bigl(\angle(\overrightarrow{OX},\overrightarrow{OY})\bigr).
\]
\end{definition}

\begin{definition}[Cross Product]
\[
\overrightarrow{OX}\times \overrightarrow{OY}=\overrightarrow{OZ}
\]
such that:
\begin{enumerate}
    \item $\overrightarrow{OX}\perp \overrightarrow{OZ}$ and $\overrightarrow{OY}\perp \overrightarrow{OZ}$, that is,
    \[
    \ip{\overrightarrow{OX}}{\overrightarrow{OZ}} = \ip{\overrightarrow{OY}}{\overrightarrow{OZ}} = 0.
    \]
    \item
    \[
    |\overrightarrow{OZ}| = |\overrightarrow{OX}|\,|\overrightarrow{OY}|\sin(\alpha),
    \]
    so $|\overrightarrow{OZ}|$ is the area of the parallelogram spanned by $\overrightarrow{OX}$ and $\overrightarrow{OY}$.
    \item If $|\overrightarrow{OZ}|\neq 0$, then the ordered set $(OX,OY,OZ)$ is right-handed.
\end{enumerate}
\end{definition}

\section{Scalar Fields}

\begin{definition}
A scalar field $\varphi$ defined on $D\subset \R^3$ is a map
\[
\varphi:D\subset \R^3\to \R.
\]
\end{definition}

\section{Vector Fields}

A vector field $F$ defined on $D\subset \R^3$ associates to each point $P\in D$ a vector $F(P)$.

With basis $\{e_1(P),e_2(P),e_3(P)\}$, the field represents the vector $F$ at $P$ as
\[
F(P)=F_1(P)e_1(P)+F_2(P)e_2(P)+F_3(P)e_3(P).
\]

In $\R^3$, a vector field is a map $F:\R^3\to \R^3$ such that for each point $(x,y,z)\in \R^3$,
\[
F(x,y,z)=\bigl(F_1(x,y,z),F_2(x,y,z),F_3(x,y,z)\bigr)\in \R^3.
\]

\section{Field Lines}

\begin{definition}
Let $F$ be a vector field on $D$. The field lines of $F$ are smooth curves $K$ such that for every $p\in K$, the vector $F(p)$ is tangent to $K$.
\end{definition}

The parametrized curve
\[
r(t)=x(t)i+y(t)j+z(t)k
\]
is a field line for the vector field $F$ if
\[
\dot r(t)=\frac{dr(t)}{dt}=\lambda(t)\,F(r(t)),
\qquad \lambda(t)>0.
\]

In polar coordinates,
\[
F(r,\theta)=F_r(r,\theta)\,\hat r + F_\theta(r,\theta)\,\hat\theta,
\]
with field line given by the system
\[
\frac{dr(t)}{dt}=\lambda(t)\,F_r\bigl(r(t),\theta(t)\bigr)
\]
and
\[
r(t)\frac{d\theta(t)}{dt}=\lambda(t)\,F_\theta\bigl(r(t),\theta(t)\bigr).
\]

\section{Potentials and Conservative Vector Fields}

If
\[
F(x,y,z)=\grad\varphi(x,y,z)\equiv \left(\frac{\partial\varphi}{\partial x},\frac{\partial\varphi}{\partial y},\frac{\partial\varphi}{\partial z}\right)
\]
on a domain $D$, then $F$ is conservative on $D$. Also, the scalar field $\varphi$ is the potential of $F$ on $D$.

\chapter{Curves, Line Integrals, and Surfaces}

\section{Curves and Arc Length}

\begin{definition}
A curve $C$ from $a$ to $b$ is the set of points traced out by a smooth path $r:[\alpha,\beta]\to \R^n$ such that:
\begin{enumerate}
    \item $r\in C^{k}$,
    \item $r(\alpha)=a$ and $r(\beta)=b$,
    \item $r$ is a bijection,
    \item $r'(t)\neq 0$ for all $t\in(\alpha,\beta)$.
\end{enumerate}
Then $r$ is a parametrization of the curve $C$.
\end{definition}

\begin{remark}
$k$ is generally $\infty$.
\end{remark}

\begin{definition}[Piecewise Smoothness]
A piecewise smooth curve is a path made of a finite number of smooth segments.
\end{definition}

\begin{note}
A closed curve is a curve with parametrization $r:[a,b]\to\R^n$ such that $r(a)=r(b)$, so it is not injective on the whole interval.
\end{note}

Let $P$ be a partition of $[a,b]$ and let $r(t)$ be a parametrization. We take
\[
P:\ a=t_0<t_1<\dots<t_N=b.
\]
Then
\[
L_P=\sum_{k=0}^{N-1}|r(t_{k+1})-r(t_k)|,
\]
and the arc length of a smooth curve $C$ is given by
\[
L(C)=\sup\{L_P\mid P \text{ is a partition of } [a,b]\}=\lim_{|P|\to 0}L_P.
\]

\begin{theorem}
Let $C$ be a smooth curve with smooth parametrization
\[
r(t)=\bigl(r_1(t),\dots,r_n(t)\bigr),\qquad t\in[a,b].
\]
Then
\[
L(C)=\int_a^b \sqrt{r_1'(t)^2+r_2'(t)^2+\cdots+r_n'(t)^2}\,dt
=\int_a^b |r'(t)|\,dt.
\]
\end{theorem}

Let $r:[a,b]\to \R^n$ be a smooth parametrization of $C$ and let
\[
a=t_0<\dots<t_N=b
\]
be a partition. Approximating $C$ by segments yields
\[
\sum_{i=0}^{N-1} f(r(t_i))\,|r(t_{i+1})-r(t_i)|.
\]
By the Mean Value Theorem,
\[
|r(t_{i+1})-r(t_i)|=|r'(\tau_i)|(t_{i+1}-t_i),
\qquad \tau_i\in(t_i,t_{i+1}).
\]
Hence the sum becomes
\[
\sum f(r(t_i))|r'(\tau_i)|\,\Delta t_i,
\]
so
\[
\lim_{\Delta t_i\to 0}\sum f(r(t_i))|r'(\tau_i)|\,\Delta t_i
\]
is a Riemann sum that converges to
\[
\int_a^b f(r(t))|r'(t)|\,dt.
\]

\begin{theorem}
The line integral along the curve $C$ is
\[
\int_C f\,ds
=\int_a^b f(r(t))|r'(t)|\,dt.
\]
\end{theorem}

\section{Work, Circulation, and Path Independence}

\subsection{Line Integral of a Vector Field (Work)}

Let $C$ be a piecewise smooth curve, let $F$ be a continuous vector field, and let
\[
\ut:=\frac{r'(t)}{|r'(t)|}
\]
be the unit tangent vector, where $r(t)$ is a parametrization of $C$. Then $\ip{F}{\ut}$ exists and
\[
\int_C \ip{F}{\ut}\,ds
\]
is the work done by $F$ along $C$.

\begin{note}
Since $ds=|r'(t)|dt$ and $\ut=\dfrac{r'(t)}{|r'(t)|}$, then
\[
\int_C \ip{F}{\ut}\,ds = \int_a^b \ip{F(r(t))}{r'(t)}\,dt,
\]
where $r(a)$ and $r(b)$ are the endpoints of $C$.
\end{note}

\begin{definition}
If $r(a)=r(b)$, then $C$ is a closed curve. If intersections may be present, then we have the cyclic integral
\[
\oint_C f\,ds = \oint_C F\cdot dr,
\]
the circulation of $F$ around $C$.
\end{definition}

\subsection{Path Independence}

In a conservative vector field, the work done is independent of path.

\begin{theorem}
Let $F$ be conservative with potential $\varphi$. Then
\[
\int_C \ip{F}{\ut}\,ds
=\int_C \ip{\grad\varphi}{\ut}\,ds
=\int_a^b \ip{\grad\varphi(r(t))}{r'(t)}\,dt
=\varphi(r(b))-\varphi(r(a)).
\]
\end{theorem}

\begin{note}
Then
\[
\oint_C \ip{F}{\ut}\,ds = \varphi(r(b))-\varphi(r(a))=0.
\]
\end{note}

\section{Surface Integrals}

\begin{note}
$\dfrac{\partial x}{\partial u}$ and $\dfrac{\partial x}{\partial v}$ span the tangent plane.
\end{note}

\begin{definition}
A surface $S\subset \R^3$ is called smooth if there exist $G\subset \R^2$ and a map $x:G\to \R^3$, continuously differentiable and bijective, such that
\[
\frac{\partial x}{\partial u}\times \frac{\partial x}{\partial v}\neq 0
\qquad \forall (u,v)\in G,
\]
so $S=x(G)$.
\end{definition}

\begin{theorem}
Let $f$ be a continuous scalar field. Then
\[
\iint_S f\,d\sigma = \iint_G f\bigl(x(u,v)\bigr)\left|\frac{\partial x}{\partial u}\times \frac{\partial x}{\partial v}\right|\,du\,dv.
\]
\end{theorem}

\begin{intuition}
$d\sigma$ is the surface-area element on $S$. When we change variables to the flat domain $G$, we still want to integrate with respect to the surface area of $S$. Therefore the values of $f$ must be taken at points on the surface, hence $f(x(u,v))$, and
\[
\left|\frac{\partial x}{\partial u}\times \frac{\partial x}{\partial v}\right|
\]
accounts for how area in $G$ scales to surface area on $S$.
\end{intuition}

\section{Orientation}

The unit normal $n(P)$ at the point $P=x(u_p,v_p)$ is
\[
n(P):=\pm \frac{\dfrac{\partial x}{\partial u}\times \dfrac{\partial x}{\partial v}}{\left|\dfrac{\partial x}{\partial u}\times \dfrac{\partial x}{\partial v}\right|}(u_p,v_p).
\]

The sign is determined by checking whether $n$ and
\[
\frac{\partial x}{\partial u}\times \frac{\partial x}{\partial v}
\]
point in the same direction:
\[
\ip{n}{\dfrac{\partial x}{\partial u}\times \dfrac{\partial x}{\partial v}}>0
\iff \text{positive orientation},
\]
and
\[
\ip{n}{\dfrac{\partial x}{\partial u}\times \dfrac{\partial x}{\partial v}}<0
\iff \text{negative orientation}.
\]

\section{Flux}

\begin{definition}
The flux of $F$ through $S$ is
\[
\iint_S \ip{F}{n}\,d\sigma.
\]
\end{definition}

It follows that
\[
\iint_S \ip{F}{n}\,d\sigma
=\pm \iint_G \left(F\bigl(x(u,v)\bigr),\frac{\partial x}{\partial u}\times \frac{\partial x}{\partial v}\right)\,du\,dv.
\]

\chapter{Differential Operators and Integral Theorems}

\section{Gradient, Divergence, and Curl}

Gradient:
\[
\grad\varphi = \frac{\partial\varphi}{\partial x}i+\frac{\partial\varphi}{\partial y}j+\frac{\partial\varphi}{\partial z}k.
\]

Jacobian:
\[
DF=
\begin{pmatrix}
\dfrac{\partial F_1}{\partial x_1} & \cdots & \dfrac{\partial F_1}{\partial x_n}\\[0.4em]
\vdots & & \vdots\\[0.4em]
\dfrac{\partial F_n}{\partial x_1} & \cdots & \dfrac{\partial F_n}{\partial x_n}
\end{pmatrix}.
\]

Divergence:
\[
\divergence F = \operatorname{trace}(DF).
\]

If $\divergence F=0$, then $F$ is solenoidal.

Curl:
\[
\curl F = \left(\frac{\partial F_3}{\partial y}-\frac{\partial F_2}{\partial z}\right)i
+ \left(\frac{\partial F_1}{\partial z}-\frac{\partial F_3}{\partial x}\right)j
+ \left(\frac{\partial F_2}{\partial x}-\frac{\partial F_1}{\partial y}\right)k.
\]
Note that $\curl F=\grad\times F$.

\begin{theorem}
If a vector field has $\curl F=0$, it is irrotational.
\end{theorem}

\section{Laplacian Operator of a Scalar Field}

\[
\Delta:=\grad \cdot \grad,
\qquad
\Delta\varphi=\divergence(\grad\varphi).
\]

The scalar field $\varphi$ is harmonic on $D$ if $\Delta\varphi=0$.
\newpage
\section{Identities}

\[
\grad(\varphi\psi)=\varphi\grad\psi+\psi\grad\varphi
\qquad
\grad\times(\grad\varphi)=0
\]
\[
\divergence(\varphi F)=\varphi\,\divergence F+(\grad\varphi)\cdot F
\qquad
\divergence(\curl F)=0
\]
\[
\grad\times(\varphi F)=\varphi(\grad\times F)+\grad\varphi\times F
\]
\[
\divergence(F\times G)=\ip{G}{\grad\times F}-\ip{F}{\grad\times G}
\qquad \text{with } (F\cdot\grad)G:=DG[F]
\]
\[
\grad\times(\grad\times F)=\grad(\grad\cdot F)-(\grad\cdot\grad)F
\]
\[
\grad\times(F\times G)=F(\grad\cdot G)-G(\grad\cdot F)+(G\cdot\grad)F-(F\cdot\grad)G
\]
\[
\grad(F\cdot G)=F\times(\grad\times G)+G\times(\grad\times F)+(G\cdot\grad)F+(F\cdot\grad)G
\]

\section{Integral Theorems}

Recall that
\[
\ut=\frac{r'(t)}{|r'(t)|}
\qquad \text{and} \qquad
ds=|r'(t)|\,dt.
\]

\subsection{Gauss' Theorem (Divergence Theorem)}

Let $R\subset \R^3$ be a region with piecewise smooth boundary $\partial R$, and let $n$ be the outward unit normal. Let $F$ and $\varphi$ be continuously differentiable on a domain containing $R$. Then
\[
\iiint_R \grad\varphi\,dV = \oiint_{\partial R} \varphi n\,d\sigma,
\]
which is the vector form.

Also,
\[
\iiint_R \divergence(F)\,dV = \oiint_{\partial R} \ip{F}{n}\,d\sigma.
\]

\subsection{Stokes' Theorem}

For $\varphi,F\in C^1(D)$, it holds that
\[
\iint_S n\times \grad\varphi\,d\sigma = \oint_{\partial S} \varphi \ut\,ds
\]
and
\[
\iint_S \ip{\curl F}{n}\,d\sigma = \oint_{\partial S} \ip{F}{\ut}\,ds.
\]

\subsection{Green's Theorem}

This is a special case of Stokes' theorem for $n=(0,0,1)$.

Let $D\subset \R^2$ be open and connected, and let $R\subset D$ be a domain with piecewise smooth positively oriented boundary $\partial R$. Let $F:D\to \R^2$. Then
\[
\iint_R\left(\frac{\partial F_2}{\partial x}-\frac{\partial F_1}{\partial y}\right)dA
= \oint_{\partial R} F_1(x,y)\,dx + F_2(x,y)\,dy
= \oint_{\partial R}\ip{F}{\ut}\,ds.
\]

\section{Vector Potential}

A vector field $F$ has a vector potential if
\[
F=\curl G.
\]

A continuous vector field $F$ is source free if
\[
\oiint_S \ip{F}{n}\,d\sigma=0,
\]
so there is zero flux for closed $S$ in $D$.

If $F$ is solenoidal, that is, if $\divergence(F)=0$, then $F$ is source free.

\begin{theorem}[Helmholtz Decomposition]
For any $F:\R^3\to \R^3$ of class $C^2$, we can decompose it into
\[
-\grad\phi+\grad\times A,
\]
that is, an irrotational part plus a solenoidal part.
\end{theorem}

\begin{note}
If $\divergence F=0$ then
\[
G:=\int_0^1 t(F(tx)\times x)\,dt
\]
is a vector potential.
\end{note}

To fulfill hypothesis $\mathcal{H}$, a subset $D\subset \R^3$ needs:
\begin{itemize}
    \item For any closed surface $S\subset D$, the domain enclosed by $S$ is completely contained in $D$.
    \item $D$ is non-periphractic, so there are no holes in $D$.
\end{itemize}

\part{Modern Vector Analysis}

\chapter{Manifolds}

\section{Topological Manifolds}

A topological manifold is a space $M$ such that for every $p\in M$, there exists an open neighborhood $U\subset M$ containing $p$ and a homeomorphism
\[
\varphi:U\to V\subset \R^n,
\]
where $V$ is open.

A space that globally is weird but locally Euclidean.

The pair $(U,\varphi)$ is a chart, and it is centered if there exists $p\in U$ such that $\varphi(p)=0$.

\section{Smooth Atlases}

This is a collection
\[
\mathcal A=\{(U_\alpha,\varphi_\alpha)\}
\]
with $\bigcup_\alpha U_\alpha=M$, where $M$ is an $n$-dimensional manifold. Here, any two charts in $\mathcal A$ are smoothly compatible.

Two charts are smoothly compatible if $U_\alpha\cap U_\beta=\varnothing$ or if
\[
\varphi_\beta\circ \varphi_\alpha^{-1}:\ \varphi_\alpha(U_\alpha\cap U_\beta)\to \varphi_\beta(U_\alpha\cap U_\beta)
\]
and
\[
\varphi_\alpha\circ \varphi_\beta^{-1}:\ \varphi_\beta(U_\alpha\cap U_\beta)\to \varphi_\alpha(U_\alpha\cap U_\beta)
\]
are smooth.

\begin{note}
Smoothly compatible is also $C^\infty$-compatible.
\end{note}

\section{Maximal Atlases}

A maximal atlas $\mathcal A$ contains all charts on $M$ that are smoothly compatible.

\section{Smooth Manifolds}

A smooth manifold is a pair $(M,\mathcal A)$ having a differential structure defined by $\mathcal A$.

A scalar function $f:M\to \R$ is smooth at $p\in M$ if in a chart,
\[
f\circ \varphi^{-1}:\varphi(U)\to \R
\]
is smooth at $\varphi(p)$.

\begin{intuition}
A differential structure is the maximal collection of charts that are smoothly compatible on $M$ and that allows us to build a local pipeline
\[
\R^n \xrightarrow{\varphi^{-1}} U\subset M \xrightarrow{f} \R
\]
so that we can perform calculus independently of chart choice.
\end{intuition}

\section{Oriented Manifolds}

A smooth manifold $M$ is oriented if its maximal atlas
\[
\mathcal A=\{(U_\alpha,\varphi_\alpha)\}
\]
is such that the Jacobian determinant of each transition map $\varphi_\beta\circ\varphi_\alpha^{-1}$ is positive on overlaps.

\section{Submanifolds}

A subset $K\subset M$ is called a submanifold of dimension $k\le m$ if, for every point $p\in K$, there exists a chart $(U,\varphi)$ with $p\in U\subset M$ such that
\[
\varphi(U\cap K)=\varphi(U)\cap \iota(\R^k)
=\{(x^1,\dots,x^m)\in \varphi(U)\subset \R^m\mid x^{k+1}=x^{k+2}=\cdots=x^m=0\},
\]
where $\iota:\R^k\to \R^m$ is the inclusion map.

\begin{note}
A submanifold is a manifold.
\end{note}

\section{Manifolds with Boundary}

An $n$-dimensional manifold $M$ with boundary is a space such that every point $p\in M$ admits a chart into either $\R^n$ or the half-space $H^n$, with
\[
H^n=\{(x_1,x_2,\dots,x_n)\in \R^n\mid x_n\ge 0\}.
\]
Transition maps between these charts are smooth.

Define the boundary by
\[
\partial M:=\{p\in M:\varphi(p)\in\{x_n=0\}\}.
\]

Define the interior by
\[
M^\circ:=\{p\in M:\varphi(p)\in\{x_n>0\}\}.
\]

\begin{theorem}
$\partial M$ is a boundary-less $(n-1)$-manifold acting as the boundary of $M$.
\end{theorem}

\begin{example}
Let $M:=T^2\times [0,1]$, where $T^2$ is a torus $S^1\times S^1$.

The local coordinates for $T^2$ are $(\theta,\phi)$ and $t$ for $[0,1]$, therefore the coordinates for $M$ are $(\theta,\phi,t)\in[0,2\pi)\times[0,2\pi)\times[0,1]$.

The boundary $\partial M$ is then given by
\[
(T^2\times\{0\})\cup(T^2\times\{1\}).
\]
The interior $M^\circ$ is therefore $T^2\times(0,1)$.
\end{example}

\chapter{Tangent and Cotangent Structures}

\section{Tangent Space}

This is the set of all possible derivations at $p$:
\[
T_pM=\{v_p:C^\infty(M)\to \R\mid v_p \text{ is a derivation at } p\}.
\]

Here $v_p$ satisfies the linearity and Leibniz rules:
\begin{itemize}
    \item $v_p[\alpha f+\beta g]=\alpha v_p[f]+\beta v_p[g]$,
    \item $v_p[fg]=f(p)\,v_p[g]+g(p)\,v_p[f]$.
\end{itemize}

In local coordinates $(x^1,\dots,x^n)$, every tangent vector can be written as
\[
v_p=v^i\frac{\partial}{\partial x^i}\Big|_p.
\]

A basis for $T_pM$ with local coordinates $(x^1,x^2,x^3,\dots,x^n)$ is
\[
T_pM=\Span\left\{\frac{\partial}{\partial x^1}\Big|_p,\dots,\frac{\partial}{\partial x^n}\Big|_p\right\}.
\]

\begin{note}
These are not powers; they are indices.
\end{note}

\section{Tangent Bundle}

Defined as
\[
TM=\bigsqcup_{p\in M} T_pM.
\]

All elements in $TM$ have the form $(p,v_p)$, where $p\in M$ and $v_p\in T_pM$.

Projection: map $\pi:TM\to M$, $\pi(p,v_p)=p$.

Fiber: $\pi^{-1}(p)=T_pM$.

\begin{theorem}
If $M$ has dimension $n$, then $TM$ has dimension $2n$.
\end{theorem}

\begin{definition}
A vector field on $M$ is a smooth section of $TM$ that assigns to each point of $M$ a vector in the fiber $\pi^{-1}(p)=T_pM$.
\end{definition}

\begin{note}
$\mathfrak X(M)$ denotes the space of all vector fields on $M$.
\end{note}

\section{Dual Spaces}

The dual space $V^*$ of $V$ is the set of linear maps $V\to \R$ that measure vectors in $V$. They both have the same dimension.

If $V$ has basis $\{e_1,e_2,\ldots,e_n\}$, then we can define a dual basis $\{\mu^1,\mu^2,\ldots,\mu^n\}$ of $V^*$ by imposing that
\[
\mu^i(e_j)=\delta_j^i.
\]

Elements in $V^*$ are covectors, and applying $\alpha\in V^*$ to $v\in V$ is the canonical pairing:
\[
\langle \alpha, v\rangle:=\alpha(v).
\]

\[
\mu^i(v)=\langle\mu^i,v\rangle = \left\langle\mu^i,v^j e_j\right\rangle = v^j\langle\mu^i,e_j\rangle = v^j\delta_j^i = v^i.
\]

\[
\delta_j^i=
\begin{cases}
1 & \text{if } i=j,\\
0 & \text{if } i\neq j.
\end{cases}
\]

\section{Cotangent Space}

$T_p^*M$ is the dual space of $T_pM$, where its elements are covectors. It is spanned by
\[
\left\{dx^1,dx^2,\dots,dx^n\right\}.
\]

\[
T^*M=\bigsqcup_{p\in M} T_p^*M,
\]
exactly as in the tangent bundle, with elements $(p,\xi)$ where $p\in M$ and $\xi\in T_p^*M$. The projection is $\pi^*$ and the fiber is $(\pi^*)^{-1}(p)$.

\section{Differential 1-Forms}

A differential $1$-form is a smooth section of the cotangent bundle that assigns to each point $p$ a covector $\omega_p\in T_p^*M$.

Equivalently, it acts on vector fields as
\[
\omega:\mathfrak X(M)\to C^\infty(M).
\]

Thus it assigns a linear map that measures tangent vectors at each point, from the fiber $(\pi^*)^{-1}(p)\subset T^*M$.

\begin{note}
$\Omega^1(M)$ denotes the space of all differential $1$-forms.
\end{note}

\chapter{Tensors, Metrics, and Differential Forms}

\section{Tensors}

A tensor of type $(q,r)$ is a multilinear function defined as follows: for $p\in M$,
\[
\mathcal T_p:
\underbrace{T_p^*M\times \cdots \times T_p^*M}_{q}
\times
\underbrace{T_pM\times \cdots \times T_pM}_{r}
\longrightarrow \R.
\]

So we have
\[
\mathcal T_p(\alpha^1,\dots,\alpha^q,v_1,\dots,v_r)
\equiv
\sum_{i_1=1}^n\cdots\sum_{i_q=1}^n
\sum_{j_1=1}^n\cdots\sum_{j_r=1}^n
\mathcal T^{\,i_1,\dots,i_q}_{\,j_1,\dots,j_r}
\,\alpha^1_{i_1}\cdots \alpha^q_{i_q}\,v_1^{j_1}\cdots v_r^{j_r}.
\]

The components of $\mathcal T_p$ are defined as
\[
\mathcal T^{\,i_1,\dots,i_q}_{\,j_1,\dots,j_r}
=
\mathcal T_p\left(dx^{i_1},\dots,dx^{i_q},
\frac{\partial}{\partial x^{j_1}},\dots,\frac{\partial}{\partial x^{j_r}}\right).
\]

\section{Riemannian Manifolds}

A Riemannian metric $g$ is a smooth assignment of $g_p$, an inner product (a tensor of type $(0,2)$), on $T_pM$. So we have
\[
g_p:T_pM\times T_pM\to \R.
\]

Note that $g_p$ is linear, symmetric, and positive definite. A pair $(M,g)$ is a Riemannian manifold.

\section{Ricci Calculus}

The flat operator $\flat:T_pM\to T_p^*M$ converts vectors to covectors. Note that
\[
\langle v^\flat,w\rangle = g_p(v,w),
\qquad \forall w\in T_pM.
\]

The sharp operator $\sharp:T_p^*M\to T_pM$ converts covectors to vectors. Note that
\[
g_p(\omega^\sharp,v)=\langle\omega,v\rangle,
\qquad \forall v\in T_pM.
\]

In a coordinate basis,
\[
g_p(v,w)=g_p\left(v^i\frac{\partial}{\partial x^i},w^j\frac{\partial}{\partial x^j}\right)=g_{ij}v^iw^j,
\]
with $g_{ij}$ the $(i,j)$-entry of the Gram matrix $G=(g_{ij})$ with inverse $G^{-1}=(g^{ij})$, such that
\[
g^{ik}g_{kj}=\delta_j^i.
\]

Note that
\[
(v^\flat)_j=g_{ij}v^i
\qquad \text{and} \qquad
(\omega^\sharp)^i=g^{ij}\omega_j.
\]

\section{Differential \texorpdfstring{$k$}{k}-Forms}

A differential $k$-form $\omega$ on $M$ is a smooth assignment of an alternating $k$-tensor $\omega_p$ acting on $k$ vectors:
\[
\omega_p:T_pM\times T_pM\times \cdots \times T_pM\to \R
\]
such that $\omega_p(v_1,v_2,\dots,v_k)=0$ whenever $v_1,\dots,v_k$ is a linearly dependent set in $T_pM$.

\begin{theorem}
An alternating $k$-tensor $\omega_p$ is also skew-symmetric. So
\[
\omega_p\bigl(v_{\pi(1)},v_{\pi(2)},\dots,v_{\pi(k)}\bigr)=\operatorname{sign}(\pi)\,\omega_p(v_1,\dots,v_k).
\]
Recall that $\operatorname{sign}(\pi)=1$ if $\pi$ is even and $-1$ if $\pi$ is odd.
\end{theorem}

So a differential $k$-form is a smooth section of $\Lambda^k(T^*M)$, that assigns to each point $p$ an alternating $k$-tensor. Here $\Lambda^k(T^*M)$ is the alternating $k$-tensor bundle.

For clarity, in coordinates:

$1$-form:
\[
\omega = \sum_{i=1}^n \omega_i\,dx^i,
\qquad \text{for example } \omega=F_1dx^1+F_2dx^2+F_3dx^3 \text{ in } \R^3.
\]

$2$-form:
\[
\eta = \sum_{i<j} \eta_{ij}\,dx^i\wedge dx^j.
\]

$k$-form:
\[
\alpha = \sum_{i_1<\cdots<i_k} \alpha_{i_1,\dots,i_k}\,dx^{i_1}\wedge\cdots\wedge dx^{i_k}.
\]

\begin{note}
There are $\binom{n}{k}$ linearly independent basis alternating $k$-tensors.
\end{note}

\section{Exterior Algebra}

Let $\mu\in\Omega^k(M)$ and $\nu\in\Omega^\ell(M)$. The wedge product is a $(k+\ell)$-differential form, so
\[
\mu\wedge\nu\in\Omega^{k+\ell}(M).
\]

In coordinates,
\[
\mu\wedge\nu
=
\mu_{i_1,\dots,i_k}\nu_{j_1,\dots,j_\ell}
\,dx^{i_1}\wedge\cdots\wedge dx^{i_k}\wedge dx^{j_1}\wedge\cdots\wedge dx^{j_\ell}.
\]

The wedge product has properties:
\begin{itemize}
    \item Distributive: $(\mu+\nu)\wedge\eta = \mu\wedge\eta + \nu\wedge\eta$.
    \item Associative: $(\mu\wedge\nu)\wedge\eta = \mu\wedge(\nu\wedge\eta)$.
    \item Anti-commutative: $\mu\wedge\eta = (-1)^{k\ell}(\eta\wedge\mu)$, where $\mu\in\Omega^k(M)$ and $\eta\in\Omega^\ell(M)$.
    \item Alternating: $dx^i\wedge dx^i=0$ for all $i$.
\end{itemize}

\begin{remark}
$\mu_{i_1,\dots,i_k}$ are the components of $\mu$ such that $i_1<i_2<\cdots<i_k$.
\end{remark}

\begin{note}
For a $0$-form $\alpha$, $\wedge$ is the same as ordinary multiplication: $\alpha\wedge\mu=\alpha\mu$.
\end{note}

\begin{example}
In $\R^3$, let $\nu\in\Omega^2(M)$,
\[
\nu_p = 2dx^1\wedge dx^2 - 3dx^1\wedge dx^3 + 5dx^2\wedge dx^3
= \sum_{i<j}\nu_{ij}\,dx^i\wedge dx^j.
\]
So then $\nu_{12}=2$, $\nu_{13}=-3$, and $\nu_{23}=5$.
\end{example}

\begin{theorem}
The wedge product is an antisymmetrized tensor product.
Formally,
\[
(\omega\wedge\mu)_p = \frac{(k+\ell)!}{k!\,\ell!}\,\mathcal O_a(\omega_p\otimes\mu_p),
\]
where $\omega\in\Omega^k(M)$ and $\mu\in\Omega^\ell(M)$.
\end{theorem}

$\mathcal O_a$ is the skew-symmetrization map, for which a $\binom{k}{0}$-tensor $\mathcal T_p$ on $T_pM$ is defined by
\[
\mathcal O_a(\mathcal T_p)(v_1,\dots,v_k)=\frac{1}{k!}\sum_{\pi}\sgn(\pi)\,\mathcal T_p(v_{\pi(1)},\dots,v_{\pi(k)}).
\]

\section{Exterior Derivative}

The exterior derivative $d$ is an operator that raises the degree of a differential form by $1$:
\[
d:\Omega^k(M)\to\Omega^{k+1}(M):\omega\mapsto d\omega,
\]
for all $k\le n$, where $\dim(M)=n$.

Properties:
\begin{itemize}
    \item Differential rule: if $\alpha\in\Omega^0(M)$, then $d\alpha$ is the differential, so
    \[
    d\alpha = \sum_i \frac{\partial\alpha}{\partial x^i}\,dx^i.
    \]
    \item Nilpotency: $d(d\alpha)=0$, hence $d^2=0$.
    \item Product rule: $d(\omega\wedge\mu)=d\omega\wedge\mu + (-1)^k\omega\wedge d\mu$.
    \item Distributivity: $d(\omega+\mu)=d\omega + d\mu$.
\end{itemize}

Also note that $d\omega$ is defined pointwise for all $p\in M$ by
\[
(d\omega)_p = d(\omega_p).
\]

\section{Closed and Exact Forms}

A differential form $\omega$ is closed if $d\omega=0$.

A differential $k$-form is exact if there exists $\mu\in\Omega^{k-1}(M)$ such that $\omega=d\mu$.

Note that exact $\Rightarrow$ closed, but closed $\nRightarrow$ exact.

\section{Contraction Operator}

\begin{definition}
Let $M$ be an $n$-dimensional manifold. With every vector field $v\in\mathfrak X(M)$ we can associate a linear map $\iota_v$, a contraction. This $\iota_v$ maps a differential $k$-form to a $(k-1)$-form, so
\[
\iota_v:\Omega^k(M)\to\Omega^{k-1}(M):\omega\mapsto \iota_v\omega,
\]
defined by
\[
(\iota_v\omega)(v_1,\dots,v_{k-1}) = \omega(v,v_1,\dots,v_{k-1})
\]
for all $v_1,v_2,\dots,v_{k-1}\in\mathfrak X(M)$.
\end{definition}

If $k=0$, then the target would be $\Omega^{-1}(M)$, so we set $\iota_vf=0$.

Properties:

Antisymmetry:
\[
\iota_v\circ \iota_w = - \iota_w\circ \iota_v.
\]

Product rule:
\[
\iota_v(\mu\wedge\omega)=(\iota_v\mu)\wedge\omega + (-1)^\ell\mu\wedge(\iota_v\omega),
\]
where $\mu\in\Omega^\ell(M)$.

\section{Push-Forward}

\begin{definition}
Let $f:M\to N$ be smooth. The push-forward (differential) of $f$ at $p\in M$ is
\[
f_*: \mathfrak X(M)\to \mathfrak X(N): v \mapsto f_*v,
\]
defined for $v\in \mathfrak X(M)$ and $g\in C^\infty(N)$ by
\[
(f_{*}v)(g)\equiv v(g\circ f).
\]
\end{definition}

\section{Pullback}

\begin{definition}
Let $f:M\to N$ be smooth. The pullback map is the linear map
\[
f^*:\Omega^k(N)\to\Omega^k(M):\omega\mapsto f^*\omega,
\]
defined for $\omega\in\Omega^k(N)$ and $v_1,\dots,v_k\in \mathfrak X(M)$ by
\[
(f^*\omega)(v_1,\dots,v_k)\equiv \omega(f_{*}v_1,\dots,f_{*}v_k).
\]
\end{definition}

Pullback map properties:
\begin{itemize}
    \item Compatible with the wedge product: $f^*(\mu\wedge\omega)=f^*\mu\wedge f^*\omega$.
    \item Compatible with the exterior derivative: $f^*(d\omega)=d(f^*\omega)$.
\end{itemize}

\chapter{Integration of Differential Forms}

\section{Integration in \texorpdfstring{$\R^3$}{R3}}

\begin{definition}
Let $\omega\in\Omega^1(D)$, let $C\subset D$ be smooth and orientable with parametrization $x(t)=\bigl(x(t),y(t),z(t)\bigr)$, $t\in[a,b]$. Then the integral of $\omega$ along $C$ is
\[
\int_C \omega \equiv \int_a^b \omega_{x(t)}(\dot x(t))\,dt.
\]
\end{definition}

It follows that by definition of $\omega$,
\[
\int_a^b \omega_{x(t)}(\dot x(t))\,dt
=\int_a^b \sum_{i=1}^3 F_i(x(t))\,dx^i(\dot x(t))\,dt
=\int_a^b \sum_{i=1}^3 F_i(x(t))\dot x^i(t)\,dt
\]
\[
=\int_a^b \ip{F(x(t))}{\dot x(t)}\,dt
=\int_C \ip{F}{\ut}\,ds.
\]

\begin{definition}
Let $\omega\in\Omega^2(D)$ and let $S\subset D$ be a smooth orientable surface with parametrization $x=x(u,v)$, $(u,v)\in G\subset \R^2$. Then the integral over $S$ is defined by
\[
\int_S \omega \equiv \iint_G \omega_{x(u,v)}\!\left(\frac{\partial x}{\partial u},\frac{\partial x}{\partial v}\right)\,du\,dv.
\]
\end{definition}

It follows that
\[
\iint_G \omega_{x(u,v)}\!\left(\frac{\partial x}{\partial u},\frac{\partial x}{\partial v}\right)\,du\,dv
=\iint_G F_1(x(u,v))\,dy\wedge dz\!\left(\frac{\partial x}{\partial u},\frac{\partial x}{\partial v}\right)\,du\,dv
\]
\[
\qquad + \iint_G F_2(x(u,v))\,dz\wedge dx\!\left(\frac{\partial x}{\partial u},\frac{\partial x}{\partial v}\right)\,du\,dv
\]
\[
\qquad + \iint_G F_3(x(u,v))\,dx\wedge dy\!\left(\frac{\partial x}{\partial u},\frac{\partial x}{\partial v}\right)\,du\,dv
\]
\[
= \iint_G \left(F(x(u,v)),\frac{\partial x}{\partial u}\times \frac{\partial x}{\partial v}\right)\,du\,dv
= \iint_S \ip{F}{n}\,d\sigma.
\]

\subsection{1-form \texorpdfstring{$\to$}{->} curl (line \texorpdfstring{$\to$}{->} surface)}

\[
\omega = F_1\,dx + F_2\,dy + F_3\,dz.
\]

\[
d\omega =
(\partial_y F_3 - \partial_z F_2)\,dy\wedge dz
+
(\partial_z F_1 - \partial_x F_3)\,dz\wedge dx
+
(\partial_x F_2 - \partial_y F_1)\,dx\wedge dy.
\]

\textbf{Identification:}
\[
d\omega \equiv (\curl F)_1\,dy\wedge dz + (\curl F)_2\,dz\wedge dx + (\curl F)_3\,dx\wedge dy.
\]

\textbf{Integral form:}
\[
\int_C \omega
=
\int_C F_i\,dx^i
=
\int_C \ip{F}{\ut}\,ds,
\qquad
\int_S d\omega
=
\iint_S \ip{\curl F}{n}\,d\sigma.
\]

\subsection{2-form \texorpdfstring{$\to$}{->} divergence (surface \texorpdfstring{$\to$}{->} volume)}

\[
\omega = F_1\,dy\wedge dz + F_2\,dz\wedge dx + F_3\,dx\wedge dy
\]
implies
\[
d\omega =
(\partial_x F_1 + \partial_y F_2 + \partial_z F_3)\,dx\wedge dy\wedge dz.
\]

\textbf{Identification:}
\[
d\omega \equiv (\divergence F)\,dx\wedge dy\wedge dz.
\]

\textbf{Integral form:}
\[
\int_S \omega
=
\iint_S \ip{F}{n}\,d\sigma,
\quad 
\int_D d\omega
=
\iiint_D \divergence F\,dV.
\]

\subsection{3-form (volume)}

\[
\omega = f\,dx\wedge dy\wedge dz
\qquad \Rightarrow \qquad
\int_D \omega = \iiint_D f\,dV.
\]

\section{Integration in \texorpdfstring{$\R^n$}{Rn}}

\subsection{Parametrized \texorpdfstring{$k$}{k}-Subsets}

A set $M\subset D\subset \R^n$ is a parametrized $k$-subset if there exists
\[
x:G\subset \R^k\to \R^n
\]
such that $x$ is injective and $M=x(G)$.

It is smooth if $x\in C^\infty$ and if the vectors
\[
\frac{\partial x}{\partial u_1},\frac{\partial x}{\partial u_2},\dots,\frac{\partial x}{\partial u_k}
\]
are linearly independent.

\subsection{Orientability}

Let $M$ be a smooth parametrized $k$-subset with parametrization
\[
x:G\subset \R^k\to \R^n.
\]
$M$ is orientable if there exists a non-zero differential $k$-form $\beta$ on $M$ such that the parametrization is compatible:
\[
\beta_x(T_{u_1},T_{u_2},\cdots,T_{u_k})>0
\qquad \forall (u_1,u_2,\cdots,u_k)\in G.
\]

\section{Integration on Manifolds}

\subsection{Integration of Top-Forms in a Single Chart}

Let $M$ be an oriented $n$-dimensional smooth manifold, and let $\omega \in \Omega^n(M)$ have compact support contained in a chart $(U,\varphi)\subset M$.

We define the integral of $\omega$ over $M$ by pulling it back to $\R^n$:
\[
\int_M \omega := \int_{\varphi(U)} (\varphi^{-1})^*\omega.
\]

In coordinates, this becomes
\[
(\varphi^{-1})^*\omega = f(x)\,dx^1 \wedge \cdots \wedge dx^n,
\]
so that
\[
\int_M \omega = \int_{\varphi(U)} f(x)\,dx^1 \cdots dx^n.
\]

\begin{note}
The integral is well-defined since $\varphi$ is a diffeomorphism and $\supp(\omega)\subset U$.
\end{note}

\subsection{Independence of the Chart}

Let $(U,\varphi)$ and $(V,\psi)$ be two charts such that $\supp(\omega)\subset U\cap V$.

Writing
\[
h = \psi \circ \varphi^{-1},
\]
we obtain, by change of variables,
\[
\int_{\psi(V)} g(y)\,dy^1 \cdots dy^n
=
\int_{\varphi(U)} g(h(x)) \det Dh(x)\,dx^1 \cdots dx^n.
\]

Thus the value of the integral does not depend on the chosen chart.

\subsection{Partition of Unity}

Let $M$ be a smooth manifold. A collection of smooth functions $\{\chi_\alpha\}_{\alpha\in A}\subset C^\infty(M)$ is called a partition of unity if:
\begin{itemize}
    \item $0 \le \chi_\alpha \le 1$,
    \item $\sum_{\alpha\in A} \chi_\alpha(p) = 1$ for all $p\in M$,
    \item the collection is locally finite.
\end{itemize}

The partition is subordinate to an open cover $\{U_\alpha\}$ if
\[
\supp(\chi_\alpha)\subset U_\alpha.
\]

\begin{theorem}
Every open cover of a smooth manifold admits a subordinate smooth partition of unity.
\end{theorem}

\subsection{Bump Functions}

A bump function is a smooth function with compact support that is identically equal to $1$ on a smaller region.

Such functions are used to construct partitions of unity.

\subsection{Integration Using Multiple Charts}

Let $\omega \in \Omega^n(M)$ with
\[
\supp(\omega)\subset \bigcup_{i=1}^s U_i,
\]
where each $U_i$ is a chart domain.

Let $\{\chi_i\}_{i=0}^s$ be a partition of unity subordinate to the open cover
\[
\{U_0, U_1, \dots, U_s\}, \qquad U_0 := M \setminus \supp(\omega).
\]

Then
\[
\int_M \omega
=
\int_M \left(\sum_{i=0}^s \chi_i\right)\omega
=
\sum_{i=0}^s \int_M \chi_i \omega.
\]

Since $\chi_0 \omega = 0$, this reduces to
\[
\int_M \omega = \sum_{i=1}^s \int_M \chi_i \omega.
\]

Each term satisfies
\[
\int_M \chi_i \omega
=
\int_{\varphi_i(U_i)} (\varphi_i^{-1})^*(\chi_i \omega),
\]
so the integral is computed locally in each chart.

\begin{note}
The result is independent of both the chosen charts and the partition of unity.
\end{note}

\section{Regions with Boundary via Regular Values}

Let $M$ be a smooth manifold and $f \in C^\infty(M)$ such that $0$ is a regular value.

Define the region
\[
D := \{p \in M \mid f(p) \ge 0\},
\]
its boundary
\[
\partial D := \{p \in M \mid f(p) = 0\},
\]
and its interior
\[
D^\circ := \{p \in M \mid f(p) > 0\}.
\]

\begin{note}
The boundary is given by a level set of $f$.
\end{note}

\section{Generalized Stokes' Theorem}

Let $M$ be an oriented $n$-dimensional manifold and $D \subset M$ a region with boundary $\partial D$. Let $\omega \in \Omega^{n-1}(M)$ be such that $\supp(\omega)\cap D$ is compact.

Then
\[
\int_D d\omega = \int_{\partial D} \omega.
\]

The boundary integral is interpreted as
\[
\int_{\partial D} \omega = \int_{\partial D} \iota^*\omega,
\]
where $\iota:\partial D \hookrightarrow M$ is the inclusion map.

\begin{intuition}
The exterior derivative $d\omega$ measures the infinitesimal behaviour of $\omega$ inside $D$, while the integral over $\partial D$ captures its global effect on the boundary.
\end{intuition}




\end{document}